\section{Related Work}
\label{sec:related}

This section reviews the literature with a focus on how a node learns
the state of its neighbors and what determines when it re-advertises its
own.
Table~\ref{tab:related} summarizes this for the learning-based UWSN
routers closest to our setting; the remaining subsections cover adaptive
HELLO schemes in Q-learning routers for other mobile networks,
change-driven control messages in ad hoc routing and vehicular
standards, adaptive beaconing and Trickle, and the
estimation and information-freshness literature.

\begin{table*}[t]
\centering
\caption{How learning-based UWSN routers keep neighbor state current
(``n.p.''~=~not published). Last row: DA-HELLO.}
\label{tab:related}
\scriptsize
\setlength{\tabcolsep}{3pt}
\begin{tabular}{@{}l l p{4.0cm} p{3.2cm} p{4.0cm}@{}}
\hline
Protocol & Learning & How neighbor state is disseminated & Control-message interval & Rule for re-advertising \\
\hline
QELAR \cite{hu2010qelar} & Q-learning (model-based) & Carried in data packets & none & Data-driven; no control message \\
CARMA \cite{divalerio2019carma} & Distributed RL & Header fields and overhearing; 6-byte HELLO for discovery & n.p. & Data-driven \\
DROR \cite{wang2023rldepth} & Q-learning + opportunistic & Header fields; sink broadcast on delivery & none & Data-driven \\
SAQR \cite{zhang2025saqr} & Q-learning & Carried in data packets & none & Data-driven \\
Chang \emph{et~al.} \cite{chang2019passive} & RL & Beacon from every node & One per time slot (slot length n.p.) & Fixed timer \\
MARV \cite{wang2024marv} & Multi-agent RL & Header fields plus a periodic control packet & n.p. & Fixed timer \\
CLORP \cite{liu2024clorP} & Multi-agent RL & HELLO message & n.p. & Fixed timer \\
QLAR \cite{li2025qlar} & Q-learning + EKF & Periodic HELLO with position, velocity, and $V$-value & $\tau\min_j T_{\mathrm{LET}}$, $\tau=0.3$ & Timer scaled by predicted link-expiration time under a fixed-range model \\
LQARP \cite{liu2026lqarp} & Q-learning & Header fields ($V$-value, neighbor count); position and energy from one initialization broadcast & none & Data-driven \\
\hline
DA-HELLO (this work) & Q-learning hosts (host-parametric) & HELLO on demand; one broadcast at first contact & Event-driven & Change could alter a neighbor's candidate set or order \\
\hline
\end{tabular}

\end{table*}

\subsection{Learning-Based UWSN Routing and Neighbor-State Maintenance}
Reinforcement-learning routers for underwater networks differ in their
learning rule, ranging from tabular Q-learning with an energy-aware
reward \cite{hu2010qelar} to channel-aware distributed learning
\cite{divalerio2019carma}, service-aware rewards \cite{zhang2025saqr},
void recovery \cite{wang2023rldepth}, multi-agent formulations
\cite{wang2024marv,liu2024clorP}, and rewards built on multipath-aware
link quality \cite{liu2026lqarp}.
They share one requirement: a forwarder needs its neighbors' values
(value estimates, residual energy, depth, or link quality) at decision
time.
The protocols meet it in three ways.

The first is to carry the state in data packets and learn it by
overhearing, as in QELAR, CARMA, DROR, SAQR, and LQARP.
This costs no control traffic, but a node that is not on a forwarding
path advertises nothing.
MARV makes the point explicit: besides the header fields, each node
broadcasts a control packet ``periodically'' so that neighbors learn its
state ``even if they do not have packets to forward''
\cite{wang2024marv}.
The second is a timer.
Chang \emph{et~al.} have every node beacon at the start of each time
slot \cite{chang2019passive}; MARV and CLORP likewise rely on a
periodic control message \cite{wang2024marv,liu2024clorP}.
In these designs the interval is a fixed constant, and in several of
them it is not reported
\cite{chang2019passive,wang2024marv,liu2024clorP}.
Nothing links it to how quickly the advertised state is changing.

The third adapts the timer.
QLAR predicts neighbor positions and velocities with an extended Kalman
filter and sets the HELLO interval to a fraction $\tau$ of the shortest
predicted link-expiration time among neighbors \cite{li2025qlar}.
It is the closest UWSN precedent to our work, and we run its
link-expiration-scaled interval as a baseline (Section~\ref{sec:eval}).
Its connectivity model is a fixed communication range, so link expiry
is a purely geometric event, and it remains a timer rather than a test
of whether the advertised change matters to any neighbor's decision.

Mobility enters the UWSN literature in other ways as well.
Chang \emph{et~al.} model passively drifting sensors but keep the
per-slot beacon \cite{chang2019passive}; cooperative multi-AUV
collection plans vehicle paths with deep Q-learning
\cite{wang2024lgdqn,liu2025mdqn,lu2025wper}; an AUV can be dispatched
to repair isolated sub-networks \cite{nandyala2025isolated};
acoustic--radio meshes switch media opportunistically
\cite{jiang2023hybrid}; routes, power, and link schedules have been
optimized jointly \cite{yuan2024multiobj,yuan2024crosslayer}; and
spatial-temporal graphs record neighbor relations over periodic
intervals for multi-UUV routing \cite{xu2025stgr}.
In these designs motion affects where a vehicle goes or which medium is
used; none uses it to decide when a neighbor should be told.

\subsection{Adaptive HELLO Intervals in Q-Learning and Ad Hoc Routing}
Outside UWSNs, several Q-learning FANET/MANET routers adapt the HELLO
period to speed, neighbor churn, or predicted link lifetime
\cite{liu2020qmr,wei2023qfagr,cui2022tarraq,mahmud2019eehello,han2013hello}.
These schemes set \emph{how often} to advertise from kinematics or
topology change rates; they do not test whether a given attribute change
would alter a learned ranking at a neighbor.
Therefore, rate adaptation alone does not address decision relevance for
Q-learning hosts.

\subsection{Change-Driven Control Messages}
Triggered updates appear in classical ad hoc routing and neighbor
discovery: DSDV triggered updates, TBRPF next-hop-relevant reporting,
NHDP responsive HELLOs, and ETSI CAM rules keyed on kinematics
\cite{perkins1994dsdv,rfc3684,rfc6130,etsi302637}.
The common pattern is to advertise when a link, metric, or trajectory
changes beyond a threshold.
However, the trigger is defined on topology or kinematics rather than on
whether a \emph{learned} score would change a neighbor's next-hop choice.
Therefore, a decision-relevance test for Q-learning hosts remains open.

\subsection{Adaptive Beaconing and Trickle}
Terrestrial networks adapt the beacon rate to speed or neighbor turnover
\cite{ingelrest2007tap,sadiq2013beaconing}.
Trickle doubles its interval while neighbors are consistent and resets on
inconsistency \cite{levis2004trickle,rfc6206}; RPL applies it to routing
control and resets on routing-state inconsistency \cite{rfc6550}.
These schemes answer \emph{how often} to beacon given mobility or churn,
not whether a given change would alter a learned candidate order.
We therefore use periodic timers and a tuned Trickle as baselines.

\subsection{Event-Triggered Estimation and the Value of Updates}
Event-triggered estimation transmits an update when a prediction residual
or innovation exceeds a threshold
\cite{priel2023consensus,sebastian2024eco,deng2023auto,deng2024lcsys,zhong2024tsp}.
Li \emph{et~al.} use media-specific triggers for a hybrid underwater
network \cite{li2026dse}; Kalman and Markov models have been compared for
predicting acoustic SNR from traces \cite{busacca2024snr}.
In all of these works the trigger protects an estimation error.
Age-of-information scheduling instead asks how often a source should
update so that the receiver's information stays fresh
\cite{kaul2012aoi,yates2021aoisurvey}; value-of-information and
goal-oriented communication transmit an update only when it matters to
the receiver's goal
\cite{ayan2019voi,kountouris2021semantics,li2024got}.
Our trigger is a concrete instance of this goal-oriented principle in
which the receiver's goal is its next-hop choice; it does not use an
estimation-error trigger.

\subsection{Positioning and Research Gap}
The principle ``advertise only what can change a neighbor's routing
decision'' is therefore not new: it appears in DSDV triggered updates,
TBRPF's next-hop-relevant reporting, NHDP responsive HELLOs, RPL's
inconsistency reset, QFAGR's hole broadcasts, TARRAQ's event-driven
sensing interval, and goal-oriented updating.
On the campaigned hosts, the realized DA-HELLO traffic under locked
configurations is dominated by first-contact replies, membership
(void-status) events, and optional slow refresh rather than by a
continuously firing rank test.
What remains open, and what this paper contributes, is narrower.
First, a formal definition of decision relevance for \emph{learned}
scores: a change is relevant if it alters whether a node is a candidate
of some neighbor or flips the sign of a score gap to a co-candidate
(Section~\ref{sec:prelim}), which no shortest-path or threshold rule
captures for Q-learning hosts.
Second, a local, host-parametric estimator of that test that a sender
can evaluate from the information it is allowed to hold, with the host
entering only through its candidate test and score.
Third, a cross-host evaluation of HELLO scheduling for Q-learning UWSN
routers with different control planes, against tuned adaptive
baselines that include QLAR's link-expiration timer, under one
energy-accounting rule that includes idle listening, and on a shared
acoustic medium with a contention MAC, half-duplex operation, and
heard-only neighbor discovery.
